% Week 4 Report: Dataset Fusion Part 2
% Multi-Disease AI Healthcare Risk Prediction Platform
% Author: [Your Name]
% Date: February 10, 2026

\documentclass[12pt,a4paper]{article}

% Packages
\usepackage[utf8]{inputenc}
\usepackage[margin=1in]{geometry}
\usepackage{amsmath}
\usepackage{amssymb}
\usepackage{graphicx}
\usepackage{booktabs}
\usepackage{hyperref}
\usepackage{listings}
\usepackage{xcolor}
\usepackage{caption}
\usepackage{subcaption}
\usepackage{float}
\usepackage{multirow}
\usepackage{enumitem}

% Code listing style
\lstset{
    basicstyle=\ttfamily\small,
    breaklines=true,
    frame=single,
    language=Python,
    showstringspaces=false,
    commentstyle=\color{gray},
    keywordstyle=\color{blue},
    stringstyle=\color{red}
}

% Title
\title{\textbf{Week 4 Technical Report:} \\ 
       Dataset Fusion Part 2 - Preprocessing Pipeline \\ 
       \large Multi-Disease AI Healthcare Risk Prediction Platform}
\author{[Your Name] \\ Roll No: [Your Roll Number] \\ 
        Supervisor: Dr. R. Selvi \\ 
        IIIT Sri City}
\date{February 10, 2026}

\begin{document}

\maketitle

\begin{abstract}
This report documents the preprocessing pipeline developed in Week 4 for the Multi-Disease AI Healthcare Risk Prediction Platform. We detail the complete data fusion process, including missing value handling, feature scaling, quality assurance, and train/validation/test split creation. Critically, we justify our departure from traditional imputation methods by implementing two parallel approaches: a masked version with explicit missingness indicators and a sentinel-based version, enabling comprehensive ablation studies. The final preprocessed dataset comprises 191,363 samples across four diseases, ready for multi-task deep learning and baseline model training.
\end{abstract}

\tableofcontents
\newpage

%=============================================================================
\section{Introduction}

\subsection{Overview}
Following the successful unification of heterogeneous disease datasets in Week 3, Week 4 focused on addressing the critical challenge of missing values while preparing production-ready datasets. The goal was to create robust, scientifically sound preprocessing pipelines that preserve medical validity while enabling state-of-the-art deep learning approaches.

\subsection{Problem Statement}
The unified dataset exhibited 60.46\% overall missingness due to structural reasons---different diseases require different diagnostic tests. Traditional imputation methods (mean/median, KNN, MICE) would create \textit{medically invalid data}, as they assume Missing Completely At Random (MCAR) or Missing At Random (MAR), which does not hold in our multi-disease context.

\subsection{Objectives}
\begin{enumerate}[leftmargin=*]
    \item Develop missing value handling strategies that preserve medical validity
    \item Create two parallel dataset versions for comprehensive model comparison
    \item Implement proper feature scaling while respecting missingness indicators
    \item Establish rigorous quality assurance protocols
    \item Generate stratified train/validation/test splits for reproducible experiments
\end{enumerate}

%=============================================================================
\section{Methodology}

\subsection{Dataset Combination}
\subsubsection{Input Data}
We began by combining four unified datasets created in Week 3:

\begin{table}[H]
\centering
\caption{Input Dataset Statistics}
\begin{tabular}{lrrr}
\toprule
\textbf{Dataset} & \textbf{Samples} & \textbf{Features} & \textbf{Target} \\
\midrule
Diabetes & 100,000 & 24 & diabetes\_risk \\
Heart Disease & 70,000 & 24 & heart\_disease\_risk \\
Liver Disease & 30,691 & 24 & liver\_disease\_risk \\
Kidney Disease & 2,032 & 24 & kidney\_disease\_risk \\
\midrule
\textbf{Combined} & \textbf{202,723} & \textbf{24} & \textbf{4 targets} \\
\bottomrule
\end{tabular}
\end{table}

\subsubsection{Combination Strategy}
Datasets were vertically concatenated using a multi-target schema:

\begin{itemize}
    \item \textbf{Shared Features:} 24 clinical features (age, BMI, blood pressure, etc.)
    \item \textbf{Disease Targets:} 4 binary columns (diabetes\_risk, heart\_disease\_risk, liver\_disease\_risk, kidney\_disease\_risk)
    \item \textbf{Metadata:} source\_dataset column for stratification
    \item \textbf{Target Encoding:} Each row has one non-NaN target (its disease), other targets are NaN
\end{itemize}

\textbf{Script:} \texttt{src/data/combine\_datasets.py}

%=============================================================================
\subsection{Missing Value Analysis}

\subsubsection{Missingness Patterns}
A comprehensive analysis revealed structured missingness patterns:

\begin{table}[H]
\centering
\caption{Missingness by Feature Type}
\begin{tabular}{lrr}
\toprule
\textbf{Feature Category} & \textbf{Missing \%} & \textbf{Pattern} \\
\midrule
Universal (age, gender) & 0\% & Complete \\
Cardiovascular (BP, cholesterol) & 63.4\% & Diabetes/liver patients \\
Metabolic (glucose, HbA1c) & 65.1\% & Heart/kidney patients \\
Liver markers (ALT, AST, bilirubin) & 90.4\% & Non-liver patients \\
Kidney markers (creatinine, GFR) & 99.0\% & Non-kidney patients \\
\midrule
\textbf{Overall Dataset} & \textbf{60.46\%} & \textbf{Structural} \\
\bottomrule
\end{tabular}
\end{table}

\subsubsection{Critical Insight: Structural vs Random Missingness}
Our missingness is \textbf{structural} (Missing By Design), not random:

\begin{itemize}
    \item \textbf{Cause:} Patients are tested only for relevant disease markers
    \item \textbf{Example:} Diabetes patients ($n=100,000$) \textit{never received} liver function tests
    \item \textbf{Implication:} No relationship exists between observed features and missing values
    \item \textbf{Consequence:} Traditional imputation would fabricate false medical data
\end{itemize}

\textbf{Script:} \texttt{src/data/analyze\_missing\_values.py}

%=============================================================================
\subsection{Imputation Strategy}

\subsubsection{Why Traditional Imputation Fails}

\paragraph{Standard Methods Assume Randomness:}
\begin{itemize}
    \item \textbf{Mean/Median:} Assumes missing values follow population distribution
    \item \textbf{KNN Imputation:} Assumes similar patients have similar values for all features
    \item \textbf{MICE:} Assumes conditional relationships between observed and missing features
\end{itemize}

\paragraph{Violation in Our Context:}
Consider a diabetes patient with missing ALT (liver enzyme):
\begin{itemize}
    \item KNN would find similar diabetes patients and impute their (non-existent) ALT values
    \item Result: Fabricated ALT=45 U/L when the patient \textit{was never tested}
    \item Medical invalidity: Model learns false correlations between diabetes features and imputed liver markers
\end{itemize}

\subsubsection{Our Dual-Strategy Approach}

We implement \textbf{two parallel preprocessing pipelines} for rigorous comparison:

\begin{table}[H]
\centering
\caption{Dual Imputation Strategy}
\begin{tabular}{p{3cm}p{5.5cm}p{5.5cm}}
\toprule
\textbf{Aspect} & \textbf{Masked Version} & \textbf{Sentinel Version} \\
\midrule
\textbf{Method} & Value + binary mask pairs & Sentinel value (-999) \\
\midrule
\textbf{Columns} & 51 (24 values + 22 masks + 4 targets + 1 metadata) & 29 (24 features + 4 targets + 1 metadata) \\
\midrule
\textbf{Encoding} & 
\texttt{bmi=25.3, bmi\_mask=1} (present) \newline
\texttt{bmi=0, bmi\_mask=0} (missing) & 
\texttt{bmi=25.3} (present) \newline
\texttt{bmi=-999} (sentinel) \\
\midrule
\textbf{Best For} & FT-Transformer (attention on masks) & XGBoost, Random Forest \\
\midrule
\textbf{Rationale} & Explicit missingness signal for neural attention & Tree splits naturally handle sparse sentinels \\
\bottomrule
\end{tabular}
\end{table}

\paragraph{Why Sentinel = -999 (Not -1):}
After StandardScaler normalization:
\begin{itemize}
    \item Real feature values: $\mu \pm 3\sigma$ (approximately [-3, +3])
    \item Sentinel -1: \textbf{Ambiguous} (overlaps with low normalized values)
    \item Sentinel -999: \textbf{Unambiguous} (clearly distinguishable)
\end{itemize}

\textbf{Scripts:}
\begin{itemize}
    \item \texttt{src/data/impute\_missing\_values\_masked.py}
    \item \texttt{src/data/impute\_missing\_values\_sentinel.py}
\end{itemize}

%=============================================================================
\subsection{Feature Scaling}

\subsubsection{Scaling Strategy}

\paragraph{Masked Version:}
\begin{lstlisting}[language=Python]
# Scale value columns only
value_cols = ['bmi', 'glucose', 'blood_pressure', ...]
scaler = StandardScaler()
scaled_data[value_cols] = scaler.fit_transform(data[value_cols])

# Mask columns remain binary (0/1)
# No scaling applied to masks
\end{lstlisting}

\paragraph{Sentinel Version:}
\begin{lstlisting}[language=Python]
# Exclude -999 from mean/std calculation
for col in feature_columns:
    real_values = data[col][data[col] != -999]
    scaler.fit(real_values.values.reshape(-1, 1))
    
    # Apply only to real values
    mask = data[col] != -999
    scaled_data.loc[mask, col] = scaler.transform(
        data.loc[mask, col].values.reshape(-1, 1)
    )
    
    # Preserve -999 sentinel
    scaled_data.loc[~mask, col] = -999
\end{lstlisting}

\subsubsection{Scaler Persistence}
Fitted scalers saved for inference pipeline:
\begin{itemize}
    \item \texttt{models/scaler\_masked.joblib}
    \item \texttt{models/scaler\_sentinel.joblib}
\end{itemize}

\textbf{Script:} \texttt{src/data/scale\_features.py}

%=============================================================================
\subsection{Quality Assurance \& Data Fixes}

\subsubsection{QA Validation Protocol}

Comprehensive checks implemented:

\begin{enumerate}
    \item \textbf{Duplicate Detection:} Exact row matching across all features
    \item \textbf{Range Validation:} Medical plausibility checks
    \item \textbf{NaN Verification:} No unexpected NaN after imputation
    \item \textbf{Target Distribution:} Class balance analysis
    \item \textbf{Mask Integrity:} Binary 0/1 verification (masked version)
    \item \textbf{Sentinel Preservation:} -999 retained post-scaling (sentinel version)
\end{enumerate}

\subsubsection{Issues Identified \& Resolved}

\begin{table}[H]
\centering
\caption{Data Quality Issues and Fixes}
\begin{tabular}{p{4cm}p{4cm}p{5cm}}
\toprule
\textbf{Issue} & \textbf{Original State} & \textbf{Resolution} \\
\midrule
\textbf{Target Encoding} & 
Liver: 1=disease, 2=no disease \newline
Diabetes: continuous [0.027, 0.672] & 
Liver: converted 2$\rightarrow$0 \newline
Diabetes: binarized at 0.5 threshold \\
\midrule
\textbf{Duplicates} & 
11,497 rows (5.7\%) \newline
- Liver: 11,323 (37\%) \newline
- Heart: 37 (<0.1\%) & 
Removed all exact duplicates \newline
Final: 191,363 unique rows \\
\midrule
\textbf{Extreme Values} & 
BMI up to 55 std \newline
ALT/AST up to 30$\times$ normal & 
Validated as medically plausible \newline
(super obesity, severe cirrhosis) \\
\bottomrule
\end{tabular}
\end{table}

\paragraph{Reasoning for Duplicate Removal:}
\begin{itemize}
    \item Liver dataset duplicates traced to data collection artifact (multiple hospital records)
    \item Retention would bias model toward liver patterns
    \item Final 137 near-duplicates retained (from different source datasets, medically valid)
\end{itemize}

\textbf{Scripts:}
\begin{itemize}
    \item \texttt{src/data/qa\_validation.py}
    \item \texttt{src/data/fix\_data\_issues.py}
\end{itemize}

%=============================================================================
\subsection{Train/Validation/Test Splits}

\subsubsection{Split Configuration}

\begin{table}[H]
\centering
\caption{Dataset Split Specification}
\begin{tabular}{lrrl}
\toprule
\textbf{Split} & \textbf{Ratio} & \textbf{Samples} & \textbf{Purpose} \\
\midrule
Training & 70\% & 133,953 & Model parameter learning \\
Validation & 15\% & 28,705 & Hyperparameter tuning \\
Test & 15\% & 28,705 & Final unbiased evaluation \\
\midrule
\textbf{Total} & \textbf{100\%} & \textbf{191,363} & \\
\bottomrule
\end{tabular}
\end{table}

\subsubsection{Stratification Strategy}

\textbf{Stratification by \texttt{source\_dataset}} ensures proportional disease representation:

\begin{table}[H]
\centering
\caption{Disease Distribution Across Splits}
\begin{tabular}{lrrrr}
\toprule
\textbf{Disease} & \textbf{Overall} & \textbf{Train} & \textbf{Val} & \textbf{Test} \\
\midrule
Diabetes & 52.3\% & 52.3\% & 52.3\% & 52.3\% \\
Heart & 36.6\% & 36.6\% & 36.6\% & 36.6\% \\
Liver & 10.1\% & 10.1\% & 10.1\% & 10.1\% \\
Kidney & 1.1\% & 1.1\% & 1.1\% & 1.1\% \\
\bottomrule
\end{tabular}
\end{table}

\subsubsection{Target Positive Rate Verification}

\begin{table}[H]
\centering
\caption{Positive Class Rates Across Splits}
\begin{tabular}{lrrr}
\toprule
\textbf{Target} & \textbf{Train} & \textbf{Val} & \textbf{Test} \\
\midrule
diabetes\_risk & 2.5\% & 2.6\% & 2.8\% \\
heart\_disease\_risk & 50.2\% & 49.2\% & 49.8\% \\
liver\_disease\_risk & 71.1\% & 71.3\% & 72.1\% \\
kidney\_disease\_risk & 76.2\% & 71.1\% & 73.1\% \\
\bottomrule
\end{tabular}
\end{table}

\subsubsection{Data Leakage Prevention}

\begin{itemize}
    \item \textbf{Same indices} used for both masked and sentinel versions
    \item \textbf{Scaler fitted only on training data}, applied to val/test
    \item \textbf{Random seed = 42} for reproducibility
    \item \textbf{Split indices saved:} \texttt{datasets/splits/split\_indices.joblib}
\end{itemize}

\textbf{Script:} \texttt{src/data/create\_splits.py}

%=============================================================================
\section{Results}

\subsection{Final Dataset Statistics}

\begin{table}[H]
\centering
\caption{Final Preprocessed Dataset Characteristics}
\begin{tabular}{lrr}
\toprule
\textbf{Metric} & \textbf{Masked} & \textbf{Sentinel} \\
\midrule
\textbf{Samples} & 191,363 & 191,363 \\
\textbf{Features} & 51 & 29 \\
\textbf{File Size (scaled)} & 98.9 MB & 55.4 MB \\
\midrule
\textbf{Train samples} & 133,953 & 133,953 \\
\textbf{Train file size} & 67.2 MB & 38.3 MB \\
\midrule
\textbf{Val samples} & 28,705 & 28,705 \\
\textbf{Val file size} & 14.4 MB & 8.2 MB \\
\midrule
\textbf{Test samples} & 28,705 & 28,705 \\
\textbf{Test file size} & 14.4 MB & 8.2 MB \\
\bottomrule
\end{tabular}
\end{table}

\subsection{Data Pipeline Summary}

\begin{figure}[H]
\centering
\begin{verbatim}
Raw Datasets (4)     Combined Dataset      Imputed (2 versions)
202,723 samples  →   202,723 samples   →   191,363 samples
                                             (duplicates removed)
                          ↓
                    Missing Analysis
                    60.46% structural
                          ↓
                ┌─────────┴─────────┐
                ↓                   ↓
         Masked Version      Sentinel Version
         51 columns          29 columns
         value + mask        -999 sentinel
                ↓                   ↓
         StandardScaler      StandardScaler
         (preserve masks)    (preserve -999)
                ↓                   ↓
         QA Validation       QA Validation
                ↓                   ↓
         70/15/15 Split      70/15/15 Split
         (stratified)        (stratified)
\end{verbatim}
\caption{Week 4 Data Processing Pipeline}
\end{figure}

%=============================================================================
\section{Rationale \& Future Plans}

\subsection{Why This Approach?}

\subsubsection{Medical Validity}
\begin{itemize}
    \item \textbf{No fabricated data:} Imputation would create false medical relationships
    \item \textbf{Preserves diagnostic uncertainty:} Missing = "not tested", not "unknown value"
    \item \textbf{Interpretable to clinicians:} Explicit missingness aligns with medical practice
\end{itemize}

\subsubsection{Machine Learning Optimality}
\begin{itemize}
    \item \textbf{FT-Transformer advantage:} Attention mechanism learns feature relevance from masks
    \item \textbf{Tree-based compatibility:} Sentinel values create natural decision boundaries
    \item \textbf{Multi-task learning:} Each disease trains on its relevant feature subset
\end{itemize}

\subsubsection{Research Rigor}
\begin{itemize}
    \item \textbf{Ablation study:} Empirical comparison validates design choices
    \item \textbf{Reproducibility:} Saved split indices and scaler parameters
    \item \textbf{Defense-ready:} Documented rationale for evaluation panel
\end{itemize}

%=============================================================================
\subsection{Future Plan: Week 6-7 Model Training}

\subsubsection{Week 6: Machine Learning Baselines (Sentinel Version)}

\paragraph{Models:}
\begin{itemize}
    \item Logistic Regression (interpretable baseline)
    \item Random Forest (non-linear interactions)
    \item XGBoost (gradient boosting)
    \item LightGBM (efficiency comparison)
\end{itemize}

\paragraph{Why Sentinel Version:}
\begin{itemize}
    \item Tree-based models naturally handle -999 as "missing" via split criteria
    \item Faster training (29 vs 51 features)
    \item Establishes performance ceiling for traditional ML
\end{itemize}

\paragraph{Expected Outcome:}
Target AUC > 0.85 across all four diseases

%-----------------------------------------------------------------------------
\subsubsection{Week 7: Deep Learning with FT-Transformer (Masked Version)}

\paragraph{Architecture:}
\begin{itemize}
    \item \textbf{Input:} 24 feature values + 22 binary masks
    \item \textbf{Embedding Layer:} Separate embeddings for values and masks
    \item \textbf{Attention:} Self-attention learns cross-feature dependencies
    \item \textbf{Multi-Task Heads:} 4 disease-specific output layers
\end{itemize}

\paragraph{Why Masked Version:}
\begin{itemize}
    \item \textbf{Explicit attention signal:} Mask=0 tells model "ignore this feature"
    \item \textbf{Learnable missingness patterns:} Attention weights adapt per disease
    \item \textbf{State-of-the-art:} Gorishniy et al. (2021) showed superiority on tabular data
\end{itemize}

\paragraph{Training Strategy:}
\begin{lstlisting}[language=Python]
# Masked multi-task loss
for disease in [diabetes, heart, liver, kidney]:
    # Only compute loss where labels exist
    mask = ~target[disease].isna()
    loss += criterion(pred[disease][mask], 
                     target[disease][mask])
\end{lstlisting}

\paragraph{Expected Outcome:}
Outperform baseline ML by 3-5 AUC points due to:
\begin{itemize}
    \item Non-linear feature interactions
    \item Shared representations across diseases
    \item Optimal missingness handling
\end{itemize}

%-----------------------------------------------------------------------------
\subsubsection{Week 8: NLP Integration}

\paragraph{Clinical Text Processing:}
\begin{itemize}
    \item Extract patient histories from MTS-Dialog and ChatDoctor datasets
    \item Fine-tune PubMedBERT on medical conversation understanding
    \item Generate 768-dim embeddings per patient
\end{itemize}

\paragraph{Hybrid Model Architecture:}
\begin{verbatim}
Structured Features → FT-Transformer → [768-dim]
                                            ↓
                                       Concatenate
                                            ↓
Clinical Text → PubMedBERT ───────────→ [768-dim]
                                            ↓
                                     Fusion MLP
                                            ↓
                                   4 Disease Heads
\end{verbatim}

\paragraph{Why Text Improves Predictions:}
\begin{itemize}
    \item Captures symptoms not in structured data (fatigue, pain patterns)
    \item Temporal disease progression from conversation history
    \item Lifestyle factors mentioned in dialogue
\end{itemize}

%-----------------------------------------------------------------------------
\subsubsection{Week 9: Explainability (Dual XAI System)}

\paragraph{SHAP Values (Feature Importance):}
\begin{itemize}
    \item Compute Shapley values for each prediction
    \item Visualize: Which features contributed to risk score?
    \item Clinical validation: Do important features align with medical knowledge?
\end{itemize}

\paragraph{LLM Explanations (Natural Language):}
\begin{itemize}
    \item Use Groq API (Llama-3-70B) to generate text explanations
    \item Input: "Patient has BMI=32, glucose=180, BP=140/90. Model predicted diabetes risk=0.87. Explain why."
    \item Output: "High risk due to: (1) Prediabetic glucose levels, (2) Hypertension, (3) Obese BMI..."
\end{itemize}

\paragraph{Why Dual XAI:}
\begin{itemize}
    \item \textbf{Clinicians:} Need interpretable, trustworthy predictions
    \item \textbf{Patients:} Require plain-language explanations
    \item \textbf{Regulatory:} FDA/CE marking requires explainability for medical AI
\end{itemize}

%=============================================================================
\section{Discussion}

\subsection{Comparison with Related Work}

\begin{table}[H]
\centering
\caption{Missing Value Handling: Literature Comparison}
\begin{tabular}{p{4cm}p{4.5cm}p{4.5cm}}
\toprule
\textbf{Approach} & \textbf{Prior Work} & \textbf{Our Method} \\
\midrule
\textbf{Mean/Median} & 
Breast cancer (UCI): Simple imputation for <5\% missing & 
Inapplicable (60\% structural missingness) \\
\midrule
\textbf{MICE} & 
MIMIC-III (Johnson 2016): Assumes MAR in ICU time-series & 
Violates assumptions (Not Missing At Random) \\
\midrule
\textbf{Deep Learning Auto-encoders} & 
TabNet (Arik 2021): Learned imputation & 
Computationally expensive, risk of overfitting \\
\midrule
\textbf{Masking} & 
FT-Transformer (Gorishniy 2021): Binary masks for attention & 
\textbf{Adopted with medical justification} \\
\midrule
\textbf{Sentinel} & 
XGBoost documentation: -999 for tree splits & 
\textbf{Adapted with scaling considerations} \\
\bottomrule
\end{tabular}
\end{table}

\subsection{Limitations \& Mitigations}

\begin{enumerate}
    \item \textbf{Imbalanced Targets:}
    \begin{itemize}
        \item Issue: Diabetes 2.8\% positive, Kidney 75\% positive
        \item Mitigation: Focal loss, class weights, SMOTE during training
    \end{itemize}
    
    \item \textbf{Kidney Disease Underrepresentation:}
    \begin{itemize}
        \item Issue: Only 2K kidney samples (1.1\%)
        \item Mitigation: Transfer learning from larger datasets, synthetic augmentation
    \end{itemize}
    
    \item \textbf{Evaluation on Mixed Missingness:}
    \begin{itemize}
        \item Issue: Test set has same structural missingness
        \item Mitigation: Additional validation on complete-case subset
    \end{itemize}
\end{enumerate}

%=============================================================================
\section{Conclusion}

Week 4 successfully established a rigorous preprocessing pipeline that balances medical validity with machine learning best practices. By rejecting traditional imputation in favor of explicit missingness encoding, we ensure:

\begin{enumerate}
    \item \textbf{Clinical Integrity:} No fabricated medical data
    \item \textbf{Model Optimality:} Missingness patterns exploited as informative signals
    \item \textbf{Research Rigor:} Dual-version approach enables comprehensive ablation studies
\end{enumerate}

The resulting datasets (191,363 samples, stratified 70/15/15 splits) provide a solid foundation for Weeks 6-9, where baseline ML, deep learning, NLP integration, and explainability modules will be developed and evaluated.

%=============================================================================
\section{Deliverables}

\subsection{Code Artifacts}
\begin{itemize}
    \item \texttt{src/data/combine\_datasets.py}
    \item \texttt{src/data/analyze\_missing\_values.py}
    \item \texttt{src/data/impute\_missing\_values\_masked.py}
    \item \texttt{src/data/impute\_missing\_values\_sentinel.py}
    \item \texttt{src/data/scale\_features.py}
    \item \texttt{src/data/qa\_validation.py}
    \item \texttt{src/data/fix\_data\_issues.py}
    \item \texttt{src/data/create\_splits.py}
\end{itemize}

\subsection{Data Artifacts}
\begin{itemize}
    \item \texttt{datasets/combined\_unified.csv} (191,363 rows)
    \item \texttt{datasets/combined\_unified\_masked\_scaled.csv} (98.9 MB)
    \item \texttt{datasets/combined\_unified\_sentinel\_scaled.csv} (55.4 MB)
    \item \texttt{datasets/splits/masked/\{train,val,test\}.csv}
    \item \texttt{datasets/splits/sentinel/\{train,val,test\}.csv}
    \item \texttt{datasets/splits/split\_indices.joblib}
\end{itemize}

\subsection{Model Artifacts}
\begin{itemize}
    \item \texttt{models/scaler\_masked.joblib}
    \item \texttt{models/scaler\_sentinel.joblib}
\end{itemize}

\subsection{Documentation}
\begin{itemize}
    \item \texttt{PROGRESS\_TRACKER.md} (updated to 29\% overall completion)
    \item \texttt{docs/TECHNICAL\_DOCS.md} (Sections 5.1-5.10 added)
    \item This report: \texttt{docs/thesis/week4\_dataset\_fusion\_report.tex}
\end{itemize}

%=============================================================================
\section*{References}

\begin{enumerate}[label={[\arabic*]}]
    \item Gorishniy, Y., Rubachev, I., Khrulkov, V., and Babenko, A. (2021). Revisiting Deep Learning Models for Tabular Data. \textit{NeurIPS 2021}.
    
    \item Arik, S. Ö., and Pfister, T. (2021). TabNet: Attentive Interpretable Tabular Learning. \textit{AAAI 2021}.
    
    \item Johnson, A. E., et al. (2016). MIMIC-III, a freely accessible critical care database. \textit{Scientific Data}, 3(1), 1-9.
    
    \item Little, R. J., and Rubin, D. B. (2019). \textit{Statistical Analysis with Missing Data} (3rd ed.). Wiley.
    
    \item Chen, T., and Guestrin, C. (2016). XGBoost: A Scalable Tree Boosting System. \textit{KDD 2016}.
    
    \item Van Buuren, S., and Groothuis-Oudshoorn, K. (2011). mice: Multivariate Imputation by Chained Equations in R. \textit{Journal of Statistical Software}, 45(3), 1-67.
\end{enumerate}

\end{document}
